\section{第五日第三讲}

Ben-Zvi 已经解释了：当把曲线放到 $\overline{\mathbb F}_q$ 上时，有 $\theta_{G,X}\leftrightarrow\theta_{\check G,\check X}$。

本讲我想回到 $\mathbb F$ 上曲线 $\Sigma$ 的情形，说明形如
\begin{equation}
  \langle P_X,f\rangle=\begin{cases}0,&\text{某些情形},\\ \text{某个 L-函数},&\text{其余情形}\end{cases}
  \label{eq:15-1}
\end{equation}
的公式可以写成与 Ben-Zvi 的解释、以及 TQFT 图景相一致的形式。

本讲我们假设 $M=T^*X$、$\check M=T^*\check X$ 是一对超球空间，且都已极化——尽管今天的一切只用 $M$ 与 $\check M$ 就能做。下面的讨论非常示意，我们会跳过种种假设与细节；仔细的版本（带有恰当的告诫、一切都钉死）见 BZSV 论文。

形如 (15.1) 的公式很多。在 BZSV 里，我们把它们改写成
\[
  \langle P_X,f\rangle=\sum_{x\in\check X^{\mathrm{fix}}}L(T_x).
\]
右端的求和按空间对称地编号：取值为 $0$ 的情形来自该空间上没有不动点，L-函数则来自同一空间的切空间。用 TQFT 的语言说，就是 $P_X\in A_G(\Sigma)$ 匹配于右端，而后者被看作集合 $\Loc_\Sigma$ 上的函数。

在左端，$f$ 是 Hecke 本征形式，它对应于对偶侧的一个单点——同态 $\rho:\pi_1(\Sigma)\to\check G$。于是 $\pi_1(\Sigma)$ 作用在 $\check X$ 上，我们可以取不动点。在右端，我们对这个 $\rho$ 的不动点 $x$ 求和；$T_x$ 是切空间，带有 $\pi_1$-表示。而 L-函数由上同调上的外代数几何化，即 $L(T_x)$ 等于 Frobenius 在
\[
  \sum_j(-1)^j\wedge^jH^1(\Sigma_{\overline{\mathbb F}_q},T_x)
\]
上的迹。

\begin{rem}
这不是 L-函数的通常定义，但它通过 Grothendieck--Lefschetz 迹公式（该公式把 L-函数表达为 Frobenius 的特征多项式）与之联系；并且对 $n\times n$ 矩阵 $g$ 有
\[
  \det(1-g)=\sum_{j=0}^n(-1)^j\operatorname{trace}(g\,|\,\wedge^j\mathbb C^n),
\]
而通常的定义是左端。
\end{rem}

我们有局部几何猜想：存在范畴的等价
\[
  \Shv(X_F/G_O)\longrightarrow\check M/\check G\ \text{上的拟凝聚层},
\]
其中 $\Hom$ 的导出空间上 Frobenius 的迹范畴化了内积，它来自局部 Plancherel 公式
\setcounter{equation}{2}
\begin{equation}
  \langle T_V\delta_0,T_W\delta_0\rangle=\int_{\check G_{\mathrm{compact}}}\chi_V\overline{\chi}_W\,(\text{$\check M$ 的 q-特征标})\,d\mu_{Haar}.
  \label{eq:15-3}
\end{equation}
% [Tutor Note: 原文的 (15.2) 在印刷本中未出现（节内计数器却已进），译本按印刷面保留并把计数器置为 2，使 (15.3) 的编号与原文一致。]

在几何 Langlands 对应下，我们有
\[
  \Shv(\Bun_G)\simeq\QC^!(\Loc_{\check G}),\qquad P_X\leftrightarrow L_{\check X}.
\]
现在我们想仿照局部 Plancherel 公式 (15.3)，用 $\Hom$ 把配对几何化：
\[
  \langle P_X,f\rangle=\sum_{x\in\check X^{\mathrm{fix}}}L(T_x).
\]
在几何侧，$f$ 对应于摩天层，几何猜想给出
\[
  \Hom(P_X,f)=\Hom(L,\delta_f).
\]
左端看起来没问题，但为什么 L-层的纤维几何化了 $\sum_x L(T_x)$？

L-层在 $\rho$ 之上的纤维，涉及 $\rho:\pi_1\to\check G$ 作用在 $\check X$ 上的导出不动点上的函数/形式。

\begin{exmp}
当 $\Gamma$ 作用在向量空间 $W$ 上时，
\begin{itemize}[leftmargin=1.6em]
  \item 通常的不动点是 $W^\Gamma=H^0(\Gamma,W)$，其函数环是对称代数。
  \item 导出不动点记录了所有 $i\ge1$ 的 $H^i(\Gamma,W)$；例如 $H^1(\Gamma,W)$ 给函数环贡献一个外代数。
\end{itemize}
\end{exmp}

由上一例的讨论得到
\[
  \text{L-层的纤维}=\bigoplus_{x\in\check X^{\mathrm{fix}}}\wedge^*H^1(T_x),
\]
而 $\wedge^*H^1$ 正是 L-函数的几何化。

\begin{rem}
在某些假设下，L-层的纤维在 BZSV 论文第 11.8 节被算出。
\end{rem}

下面谈几何与算术的联系。原则上，$\mathbb F_q$ 上曲线 $\Sigma$ 的几何陈述蕴含 $\mathbb F_q$ 上曲线 $\Sigma$ 的数值陈述。实践中两种设定各有长处，相互启发。几何带来了非常重要的东西：两侧的对称性。
